\documentclass[11pt,a4paper]{article}
\usepackage[T1]{fontenc}
\usepackage[utf8]{inputenc}
\usepackage{lmodern,microtype}
\usepackage[a4paper,margin=26mm,headheight=15pt]{geometry}
\usepackage{amsmath,amssymb,amsthm,mathtools,bm}
\usepackage{booktabs,array,longtable}
\usepackage{xcolor}
\definecolor{navy}{RGB}{26,52,74}
\definecolor{teal}{RGB}{0,98,108}
\usepackage[colorlinks=true,linkcolor=navy,citecolor=teal,urlcolor=teal,breaklinks=true]{hyperref}
\usepackage{bookmark,fancyhdr,enumitem}
\usepackage{listings}
\lstset{basicstyle=\ttfamily\small,breaklines=true,columns=fullflexible,frame=single}
\pagestyle{fancy}
\fancyhf{}
\fancyhead[L]{\small Endpoint crossover of critical transseries}
\fancyhead[R]{\small Uniform inversion and action budgets}
\fancyfoot[C]{\thepage}
\renewcommand{\headrulewidth}{0.3pt}
\setlength{\parskip}{4pt}
\setlength{\parindent}{0pt}
\setlength{\emergencystretch}{2em}
\numberwithin{equation}{section}
\newtheorem{theorem}{Theorem}[section]
\newtheorem{lemma}[theorem]{Lemma}
\newtheorem{proposition}[theorem]{Proposition}
\newtheorem{corollary}[theorem]{Corollary}
\theoremstyle{definition}
\newtheorem{definition}[theorem]{Definition}
\newtheorem{remark}[theorem]{Remark}
% ed. (2026-09-29): unnumbered environment for ProveIt editorial notes; it does not shift any numbering.
\theoremstyle{remark}
\newtheorem*{ednoteinner}{Editorial note (ProveIt, 2026-09-29)}
\newenvironment{ednote}{\begin{ednoteinner}\sloppy\Urlmuskip=0mu plus 1mu\relax}{\end{ednoteinner}}
\newcommand{\eps}{\varepsilon}
\newcommand{\Li}{\operatorname{Li}}
\newcommand{\E}{\mathbb E}
\newcommand{\Pp}{\mathbb P}
\newcommand{\R}{\mathbb R}
\newcommand{\C}{\mathbb C}
\newcommand{\N}{\mathbb N}
\newcommand{\ee}{\mathrm e}
\newcommand{\ii}{\mathrm i}
\newcommand{\ph}{\varphi}
\newcommand{\hc}{\mathcal H}
\newcommand{\FF}{F_{\eps}}
\newcommand{\KK}{K_{\eps}}
\newcommand{\asympK}{\asymp_{K}}
\newcommand{\repo}{\textsc{ProveIt}}
\DeclareMathOperator{\Var}{Var}
\DeclareMathOperator{\Rea}{Re}
\DeclareMathOperator{\supp}{supp}
\hypersetup{pdftitle={Through the Stable--Gaussian Endpoint: Uniform Critical Transseries, Prefix-Independent Cutoff Corrections, and Conditional Extremes},pdfauthor={Research article prepared for Vladimir Reshetnikov},pdfsubject={Joint tail-index and large-order asymptotics at alpha=2}}
\begin{document}
% ed. (2026-09-29): no page anchors on the title page, which the titlepage environment numbers 1 like the following page (removes a duplicate-destination warning).
\hypersetup{pageanchor=false}
\begin{titlepage}
\vspace*{9mm}
{\color{teal}\large RESEARCH ARTICLE\par}
\vspace{8mm}
{\color{navy}\LARGE\bfseries Through the Stable--Gaussian Endpoint\par}
\vspace{5mm}
{\Large Uniform critical transseries, prefix-independent cutoff corrections, and conditional extremes\par}
\vspace{9mm}
{\large Prepared for Vladimir Reshetnikov\par}
{29 September 2026\par}
\vspace{9mm}
\begin{abstract}
We study the critical countable-action equation
\(U(q)=\sum_{j\ge1}w_jq^j\ee^{jU(q)}\), with nonnegative weights,
\(\sum jw_j=1\), and an eventually exact tail
\(w_j=c_{\eps}j^{-3+\eps}\). The tail exponent approaches two while
\(|\eps|\log n\) remains bounded. This is the joint-limit question explicitly
left open in the companion marginal-critical manuscript. We prove a uniform,
pole-cancelled inverse representation and a convergent correction chart about
a deformed Lambert core. A direct Fourier argument gives the complete central
coefficient through its first correction and a uniform local limit theorem.
For action cutoffs we obtain a two-term formula in which all finite-prefix
constants cancel. It yields a two-term minimal-budget formula and the exact
criterion for asymptotically lossless truncation. Conditioned on the coefficient
constraint, the rescaled large actions converge to a Poisson point process.
The analysis works from both sides of the endpoint, including a triangular
finite-variance regime in which the total variance is not the correct central
normalization. Proofs are self-contained apart from explicitly identified
classical analytic identities; reproducible algebraic and numerical checks
are supplied.
\end{abstract}
\vfill
\small
\textbf{Status and scope.} The proposed contribution is the uniform joint-limit
theory and its explicit correction formulas for the stated tail class, not the
classical fixed-exponent stable or Gaussian laws. Publication priority has not
been established. The article contains conventional proofs, not a new Lean
formalization or an independently peer-reviewed result. The numerical checks
are diagnostics, not interval certificates.
\end{titlepage}
\hypersetup{pageanchor=true}
\tableofcontents
\clearpage

\section{The research question and its precise resolution}
\label{sec:intro}

\subsection{The gap in the existing program}
The inspected \repo\ transseries index describes a large canonical inversion
volume, focused Lean modules, and separate research arrivals
\cite{repo-index,repo-readme}. Among those arrivals,
\emph{Beyond Finite-Action Folds} treats positive eventually exact tails
\(j^{-1-\alpha}\), for fixed \(1<\alpha<2\), with stable coefficient laws and
critical Hahn inversion \cite{repo-stable}. The later companion
\emph{Marginal Critical Transseries} treats the endpoint \(\alpha=2\), including
Lambert charts and cutoff corrections \cite{marginal}.
Its first proposed research question asks for joint inverse and coefficient
expansions when \((2-\alpha)\log n\) remains bounded, and specifically asks
whether the prefix-independent correction has a uniform interpolating form.
That is the target of this article.
% ed. (2026-09-29): the companion is now filed, the same question is posed in a second package, and an independent second answer was filed with this package (batch 49).
\begin{ednote}
The marginal companion is now filed in this series (batch~48) as
\path{Marginal_Critical_Transseries_Action_Budgets/}. The same question is
Question~1 of the logarithmic-endpoint package filed beside it
\cite{ed:lce}, an independent treatment of the same endpoint, which this
article did not see. The confluent package \cite{ed:cct}, filed with this
one and written independently of it, answers the question at leading order
for every rate \(\eps_n\to0\) of either sign. The two agree term by term:
its \(A_n\) and \(B_n\) are \(N\) and \(B\) here, its window forms are
those of Section~\ref{sec:interpretation}, its minimal budget is the
leading term of Corollary~\ref{cor:budget} at \(\alpha=2\), its cutoff profile
\(e^{-1/(2s^2)}\) is the \(\alpha\to2\) limit of \(\exp(-\ell^{-\alpha}/\alpha)\),
and both prove the conditioned Poisson limit. Inside the window this article
goes one order further. At \(\eps=0\) its first cutoff correction is the
marginal companion's critical correction and, after a change of
logarithmic normalization, that of the logarithmic-endpoint package.
\end{ednote}

The difficulty is not solved by substituting \(\alpha=2\) in a stable formula.
Two terms, one involving \(\Gamma(-\alpha)\) and the other involving
\(\zeta(\alpha-1)\), separately diverge. Their cancellation is of leading
order in the joint limit. Keeping one and discarding the other changes the
normalization. We keep their difference before either inversion or Fourier
analysis.

\subsection{What is proved here}
Let \(\eps=2-\alpha\). Define the entire crossover function
\begin{equation}
 \hc(z)=\frac{\ee^z-1}{z},\qquad \hc(0)=1,
 \qquad \FF(H)=H\hc(\eps H).
 \label{eq:Hfunction}
\end{equation}
For coefficient order \(n\), the main scales are determined by
\begin{equation}
 B^2=nc_{\eps}\FF(\log B),\quad H=\log B,\quad
 N=(nc_{\eps})^{1/(2-\eps)},\quad
 E=\ee^{\eps H},\quad F=\FF(H),\quad r=E/F.
 \label{eq:scales-intro}
\end{equation}
The root defining \(B\) is the large positive root. Under
\(|\eps|\log n\le K\), it exists uniquely in the range relevant below.
Then \(B\) is the central fluctuation scale and \(N\) is the action-cutoff
scale. They satisfy \(N/B\asymp_K H^{-1/2}\).

Writing \(u_n=[q^n]U\), and denoting by \(u_{n,M}\) the coefficient after
removing all actions larger than \(M\) \emph{without renormalizing the retained
weights}, we prove, uniformly for \(\ell=M/N\) in compact subsets of
\((0,\infty)\),
\begin{equation}
 \frac{u_{n,M}}{u_n}
 =\exp\left(-\frac{\ell^{-\alpha}}{\alpha}\right)
 \left[
 1+\frac r4\left\{\log\frac{2}{r\ell^2}+1-\gamma-\frac{2}{\ell^2}\right\}
 +O_K\left(\frac{(\log H)^2}{H^2}\right)
 \right].
 \label{eq:main-ratio-intro}
\end{equation}
Here \(\gamma\) is Euler's constant; constants in the error may also depend on
the chosen compact \(\ell\)-interval and on the fixed weight family.
No finite-prefix moment appears in the displayed correction.
At \(\eps=0\), \(r=1/H\), and this reduces to the correction in the marginal
companion. The result here is its uniform extension through the endpoint,
with a proved remainder on both sides.

The inverse theorem, central coefficient theorem, cutoff theorem,
minimal-budget corollary, and conditional point-process theorem are proved
in Sections~\ref{sec:inverse}--\ref{sec:extremes}. None is deduced by an
unjustified interchange of the limits in \(n\) and \(\alpha\).

\subsection{Classical results and the novelty boundary}
The Lagrange coefficient mechanism is classical. The connection between
positive coefficient models, random allocations, and conditioned sums is
also classical; Janson's treatment includes both fixed-index stable local
limits and the \(\sqrt{n\log n}\) endpoint law
\cite[Theorem 17.14 and Example 18.29]{janson}.
The distinction between hard and soft truncation of heavy tails is an
established subject \cite{chakrabarty}. We do not claim these mechanisms as
new. The exact polylogarithm continuation formula is likewise classical
\cite[Eq. 25.12.12]{dlmf-polylog}.

The statements offered as a model-specific extension are the combined
\emph{uniform} expansions, the first correction and its prefix cancellation,
the corresponding convergent deformed-core chart, and the consequences for
conditional actions in this moving-index regime. The literature inspection
is targeted, not an exhaustive priority search. A proof in this article
establishes its stated mathematical conclusion; it does not establish that
no equivalent formula has appeared elsewhere.

\section{The normalized family and exact coefficient identities}
\label{sec:model}

\begin{definition}[Admissible finite-prefix family]
Fix \(\eps_0\in(0,1/4)\) and an integer \(J\ge1\). For
\(|\eps|<\eps_0\), suppose
\begin{equation}
 w_j(\eps)=c_{\eps}j^{-3+\eps}+p_j(\eps),\qquad
 p_j=0\quad(j>J),
 \label{eq:weights}
\end{equation}
where \(c_{\eps}>0\) and the finitely many \(p_j\) are real analytic near zero.
Assume \(w_j\ge0\), \(w_1\) is bounded away from zero, and
\begin{equation}
 \sum_{j\ge1}jw_j(\eps)=1.
 \label{eq:criticality}
\end{equation}
All uniform assertions are made in a sufficiently small closed parameter
interval on which these conditions hold. Individual \(p_j\) may be negative.
Set
\begin{equation}
 G_{\eps}(z)=\sum_{j\ge1}w_jz^j
 =c_{\eps}\Li_{3-\eps}(z)+P_{\eps}(z),\qquad
 v_{\eps}=\frac{1}{c_{\eps}}\sum_{j=1}^J j^2p_j,
 \label{eq:G}
\end{equation}
where \(P_{\eps}(z)=\sum_{j=1}^Jp_jz^j\).
\end{definition}

A particularly transparent test family is
\begin{equation}
 c_{\eps}=\frac{1}{\zeta(2-\eps)+\eta},\quad
 p_1=\eta c_{\eps},\quad p_j=0\ (j>1),\quad \eta\ge0.
 \label{eq:test-family}
\end{equation}
It has \(v_{\eps}=\eta\). Taking \(\eta=0\) gives the pure polylogarithm
model; taking \(\eta=2\) changes only the raw first-action weight before
critical normalization.

Let \(U_{\eps}\in q\R[[q]]\) be the solution of
\begin{equation}
 U_{\eps}(q)=G_{\eps}\bigl(q\ee^{U_{\eps}(q)}\bigr).
 \label{eq:feedback}
\end{equation}
Coefficientwise substitution determines the solution uniquely, one degree
at a time. All its coefficients are nonnegative. The cutoff solution uses
\(G_{\eps,M}(z)=\sum_{j\le M}w_jz^j\) in place of \(G_{\eps}\).
The parameter \(\eps\) is fixed during each coefficient extraction, but may
depend on \(n\) in an asymptotic statement.

\begin{proposition}[Exact coefficient--Poisson correspondence]
\label{prop:poisson}
Put \(\rho_{\eps}=\exp(-G_{\eps}(1))\). For independent random variables
\(Z_j\sim\operatorname{Poisson}(nw_j)\), define
\(S_n=\sum_{j\ge1}jZ_j\) and \(S_{n,M}=\sum_{j\le M}jZ_j\).
These sums are almost surely finite, and \(\E S_n=n\). Then
\begin{align}
 u_n&=\frac1n[z^n]\ee^{nG_{\eps}(z)}
     =\frac{\rho_{\eps}^{-n}}n\Pp(S_n=n),\label{eq:coeff-poisson}\\
 \frac{u_{n,M}}{u_n}
 &=\ee^{-n\Lambda_M}
   \frac{\Pp(S_{n,M}=n)}{\Pp(S_n=n)},
 &\Lambda_M&=\sum_{j>M}w_j.\label{eq:cutoff-poisson}
\end{align}
Equivalently, the ratio in \eqref{eq:cutoff-poisson} is
\(\Pp(Z_j=0\text{ for every }j>M\mid S_n=n)\).
\end{proposition}
\begin{proof}
Set \(W=q\ee^{U}\); then \(W=q\ee^{G(W)}\). Lagrange inversion gives
\[
 [q^n]G(W)=\frac1n[z^{n-1}]G'(z)\ee^{nG(z)}
 =\frac1{n^2}[z^{n-1}](\ee^{nG(z)})'
 =\frac1n[z^n]\ee^{nG(z)}.
\]
The sum of the Poisson rates is finite, since \(\sum w_j<\infty\).
The total number of jumps is consequently finite almost surely. Its weighted
mean is \(n\sum jw_j=n\). The probability generating function of \(S_n\)
is \(\exp(n(G(z)-G(1)))\), proving the first identity. Apply the same
argument to the cutoff and divide. Independence of the retained and deleted
Poisson variables proves the conditional interpretation.
\end{proof}

\begin{proposition}[The real critical branch]
\label{prop:branch}
For \(0<q<\rho_{\eps}\), the real branch is parameterized uniquely by
\begin{equation}
 W=\ee^{-s},\quad q=\ee^{-s-G_{\eps}(\ee^{-s})},\quad
 U=G_{\eps}(\ee^{-s}),\quad s>0.
 \label{eq:param}
\end{equation}
Its power series has radius \(\rho_{\eps}\). If
\(\Delta=\log(\rho_{\eps}/q)\), then
\begin{equation}
 \Delta=s+G_{\eps}(\ee^{-s})-G_{\eps}(1),\qquad
 U=G_{\eps}(1)+\Delta-s.
 \label{eq:delta}
\end{equation}
\end{proposition}
\begin{proof}
The function \(q(W)=W\ee^{-G(W)}\) increases strictly on \((0,1)\), since
\(WG'(W)<\sum jw_j=1\). It runs from zero to \(\rho\).
The positive formal solution converges at every such \(q\): its successive
coefficientwise approximants are bounded by this real solution. Hence its
radius is at least \(\rho\). But
\(dU/dq=G'(W)\ee^{G(W)}/(1-WG'(W))\) tends to infinity as \(W\uparrow1\),
so the radius cannot exceed \(\rho\). Taking logarithms gives
\eqref{eq:delta}.
\end{proof}

In transseries notation, put \(q=\ee^{-x}\). Then
\(U(\ee^{-x})=\sum_{n\ge1}u_n\ee^{-nx}\), convergent for
\(x>G_{\eps}(1)\). The large-order questions concern these exponential-sector
coefficients. The critical inverse concerns
\(\Delta=x-G_{\eps}(1)\downarrow0\), rather than a claim about arbitrary
transseries or arbitrary summation procedures.

\section{Cancelling the endpoint pole before taking limits}
\label{sec:kernel}

Define analytic functions of \(\eps\) near zero by removable continuation:
\begin{equation}
 g_{\eps}=2\Gamma(-2+\eps)-\frac1\eps,\qquad
 z_{\eps}=\zeta(1-\eps)+\frac1\eps,\qquad
 a_{\eps}=1+\eps g_{\eps}.
 \label{eq:regularized-constants}
\end{equation}
Their endpoint values are
\begin{equation}
 g_0=\frac32-\gamma,\qquad z_0=\gamma,\qquad a_0=1.
 \label{eq:regularized-values}
\end{equation}
We write
\begin{equation}
 \KK(L)=\FF(L)+g_{\eps}\ee^{\eps L}+z_{\eps}+v_{\eps}.
 \label{eq:exact-kernel}
\end{equation}
For \(\eps\ne0\), this is exactly
\(2\Gamma(-2+\eps)\ee^{\eps L}+\zeta(1-\eps)+v_{\eps}\),
but \eqref{eq:exact-kernel} is the stable expression at \(\eps=0\).
In particular,
\begin{equation}
 K_0(L)=L+\frac32+v_0,
 \qquad \KK'(L)=a_{\eps}\ee^{\eps L}.
 \label{eq:kernel-derivative}
\end{equation}

\begin{lemma}[Exact desingularized expansion]
\label{lem:kernel}
There is a function \(R_{\eps}(s)\), jointly analytic near
\((\eps,s)=(0,0)\), such that, with a consistent logarithm of \(s\),
\begin{equation}
 s+G_{\eps}(\ee^{-s})-G_{\eps}(1)
 =\frac{c_{\eps}s^2}{2}
   \left\{\KK(\log(1/s))+sR_{\eps}(s)\right\}.
 \label{eq:exact-delta-kernel}
\end{equation}
More explicitly,
\begin{equation}
 R_{\eps}(s)=\frac2{c_{\eps}}
 \sum_{k\ge3}\frac{(-1)^k}{k!}
 \left(c_{\eps}\zeta(3-\eps-k)+\sum_{j=1}^Jp_jj^k\right)s^{k-3}.
 \label{eq:Rseries}
\end{equation}
The series converges normally on sufficiently small product neighborhoods.
\end{lemma}
\begin{proof}
Use the classical convergent expansion, for \(|s|<2\pi\) on a logarithmic
branch and nonintegral order,
\[
 \Li_{3-\eps}(\ee^{-s})
 =\Gamma(-2+\eps)s^{2-\eps}
   +\sum_{k\ge0}\zeta(3-\eps-k)\frac{(-s)^k}{k!}.
\]
This is the specialization of \cite[Eq. 25.12.12]{dlmf-polylog}.
The constant term cancels. The linear term cancels by
\eqref{eq:criticality}. The quadratic and nonanalytic terms combine into
\(c_{\eps}s^2\KK(\log(1/s))/2\). Terms of order at least three give
\eqref{eq:Rseries}. Their zeta factors have no parameter pole near zero;
normal convergence on smaller disks follows from the same polylogarithm
expansion, or from the zeta functional equation and the geometric factor
\((|s|/(2\pi))^k\) times a polynomial in \(k\).
Thus the remaining parameter singularity is removable, proving the formula
also at zero. Finite-prefix exponentials are entire in \(s\).
\end{proof}

\begin{lemma}[Uniform logarithmic estimates]
\label{lem:log-estimates}
Fix \(K<\infty\). As \(L\to\infty\), uniformly for \(|\eps|L\le K\),
\begin{align}
 \FF(L)&\asymp_K L,& \ee^{\eps L}&\asymp_K1,
 &\KK(L)&=\FF(L)+O_K(1),\label{eq:kernel-estimates}\\
 \KK(L)&=\FF(L)+(\tfrac32-\gamma)\ee^{\eps L}
                  +\gamma+v_{\eps}+O_K(L^{-1}).\label{eq:kernel-two-term}
\end{align}
For real \(b=O(\log L)\),
\begin{equation}
 \FF(L+b)=\FF(L)+\ee^{\eps L}b
     +O_K\left(\frac{b^2}{L}\right).
 \label{eq:F-shift}
\end{equation}
\end{lemma}
\begin{proof}
The continuous positive function \(\hc\) has positive lower and finite upper
bounds on \([-K,K]\). Analyticity gives
\(g_{\eps}=g_0+O(\eps)\), \(z_{\eps}=\gamma+O(\eps)\).
Finally, \(\FF'=\ee^{\eps L}\), \(\FF''=\eps\ee^{\eps L}\);
Taylor's formula gives the last estimate. These arguments also give the
analogous complex estimates with factors
\(\ee^{|\eps b|}\) when needed below.
\end{proof}

The term \((\ee^{\eps L}-1)/\eps\), not \(\ee^{\eps L}/\eps\), is the
crossover kernel. The subtraction is indispensable when \(\eps L\) stays
bounded. It retains the cancellation that is lost by taking a fixed-index
stable approximation first.

\section{A uniform inverse and a convergent deformed-core chart}
\label{sec:inverse}

\begin{theorem}[Two-term joint inverse]
\label{thm:inverse}
Let \(\Delta\downarrow0\), put \(D=\log(1/\Delta)\), and assume
\(|\eps|D\le K\). Define
\begin{equation}
 F_*=\FF(D/2),\quad E_*=\ee^{\eps D/2},\quad
 A_* =\frac12\log\frac{c_{\eps}F_*}{2},\quad
 d_*=(\tfrac32-\gamma)E_*+\gamma+v_{\eps}.
 \label{eq:inverse-vars}
\end{equation}
The solution of \eqref{eq:delta} has the uniform expansion
\begin{equation}
 s=\sqrt{\frac{2\Delta}{c_{\eps}F_*}}
 \left[
 1-\frac{E_*A_*+d_*}{2F_*}
 +O_K\left(\frac{(\log D)^2}{D^2}\right)
 \right].
 \label{eq:inverse-two-term}
\end{equation}
Consequently,
\begin{equation}
 s=2\sqrt{\frac{\Delta}
 {c_{\eps}D\hc(\eps D/2)}}
 \left[1+O_K\left(\frac{\log D}{D}\right)\right].
 \label{eq:inverse-leading}
\end{equation}
The corresponding solution of the original feedback equation is
\(U=G_{\eps}(1)+\Delta-s\).
\end{theorem}
\begin{proof}
Write \(L=\log(1/s)\). By \eqref{eq:exact-delta-kernel},
\begin{equation}
 2L=D+\log(c_{\eps}/2)
 +\log\{\KK(L)+\ee^{-L}R_{\eps}(\ee^{-L})\}.
 \label{eq:inverse-log-eq}
\end{equation}
First bracket the solution with \(L\) equal to
\(D/2+C\log D\), using Lemma~\ref{lem:log-estimates} and two fixed constants
\(C\) of opposite sufficiently large size. Monotonicity follows from
Proposition~\ref{prop:branch}; hence \(L=D/2+O_K(\log D)\).
All the subsequent estimates are therefore within a fixed bounded
\(\eps L\)-window.

Set \(L=D/2+A_*+b\). Initially the equation gives
\(b=O_K(\log D/D)\). By \eqref{eq:F-shift} and
\eqref{eq:kernel-two-term},
\[
 \KK(L)=F_*+E_*A_*+d_*
       +O_K((\log D)^2/D).
\]
The term \(E_*b\) is within this error. The analytic remainder in
\eqref{eq:inverse-log-eq} is exponentially smaller. Expanding the logarithm
now yields
\[
 b=\frac{E_*A_*+d_*}{2F_*}
      +O_K((\log D)^2/D^2).
\]
Exponentiate \(s=\ee^{-D/2-A_*-b}\). Since
\(\ee^{-D/2-A_*}=\sqrt{2\Delta/(c_{\eps}F_*)}\), the theorem follows.
\end{proof}

The two-term formula is an asymptotic result. A different construction gives
an actually convergent correction chart.

\begin{theorem}[Uniform analytic correction about the deformed core]
\label{thm:chart}
For \(|\eps|D\le K\) and small \(\Delta>0\), let \(L_0\) be the unique
large solution of
\begin{equation}
 \Delta=\frac{c_{\eps}}2\ee^{-2L_0}\KK(L_0),\qquad s_0=\ee^{-L_0}.
 \label{eq:deformed-core}
\end{equation}
Put \(X=s_0\), \(Y=1/\KK(L_0)\), and
\(Z=\ee^{\eps L_0}/\KK(L_0)\). There is a single holomorphic germ
\(V(\eps,X,Y,Z)\), on a product neighborhood independent of the particular
joint-limit path, such that the exact solution is
\begin{equation}
 s=s_0\exp\{-V(\eps,X,Y,Z)\}.
 \label{eq:chart}
\end{equation}
The germ has a convergent multivariate power series, is divisible by \(XY\),
and satisfies
\begin{equation}
 V=\frac{XYR_{\eps}(0)}{2-a_{\eps}Z}
       +O(X^2Y)
 \label{eq:chart-first}
\end{equation}
in a smaller fixed polydisk. In particular, on the real joint-limit paths,
\(s/s_0=1+O_K(s_0/\log(1/s_0))\).
\end{theorem}
\begin{proof}
Lemma~\ref{lem:log-estimates} shows that the right side of
\eqref{eq:deformed-core} is positive and strictly decreasing at the large
root: its logarithmic derivative is
\(-2+\KK'(L)/\KK(L)=-2+O_K(1/L)\).
The same bracket used for Theorem~\ref{thm:inverse} gives existence and
uniqueness in the large-root region.

The exact identity
\[
 \KK(L_0+v)-\KK(L_0)
 =a_{\eps}\ee^{\eps L_0}v\hc(\eps v)
\]
turns \eqref{eq:exact-delta-kernel}, with \(s=s_0\ee^{-v}\), into
\begin{equation}
 \ee^{-2v}\left[
 1+a_{\eps}Zv\hc(\eps v)
 +XY\ee^{-v}R_{\eps}(X\ee^{-v})\right]=1.
 \label{eq:analytic-chart-equation}
\end{equation}
Its left side minus one is holomorphic in all five variables. At
\((\eps,X,Y,Z,v)=(0,0,0,0,0)\), the derivative with respect to \(v\) is
\(-2\). The analytic implicit-function theorem gives one germ on a fixed
polydisk, not a different germ for each asymptotic path. At \(X=0\) or
\(Y=0\), its small solution is identically zero. Analytic division therefore
gives the factor \(XY\). Taking the coefficient of \(X\) in
\eqref{eq:analytic-chart-equation} gives \eqref{eq:chart-first}; the
remainder retains the factor \(Y\). Finally,
\(X\to0\), \(Y,Z=O_K(1/L_0)\), so every sufficiently late real path lies
inside that fixed polydisk.
\end{proof}

\begin{remark}[What is and is not a finite-generator statement]
The correction \(V\) is a genuinely convergent finite-generator analytic
chart. Its core is the inverse of \eqref{eq:deformed-core}; it is not asserted
to be an elementary function, nor an ordinary Lambert function when
\(\eps\ne0\). We do not claim a global complex branch classification or a
convergent expansion uniform over fixed nonzero \(\eps\) as \(D\to\infty\).
\end{remark}

At \(\eps=0\), put \(d=3/2+v_0\). The core becomes
\(\Delta=(c_0/2)\ee^{-2L_0}(L_0+d)\), so it is explicitly
\begin{equation}
 s_0=\left[
 \frac{-4\Delta}
 {c_0 W_{-1}(-4\Delta\ee^{-2d}/c_0)}\right]^{1/2}.
 \label{eq:lambert-core}
\end{equation}
Thus Theorem~\ref{thm:chart} deforms the real negative Lambert branch
continuously and analytically at the level of the desingularized equation.
The identity follows directly by setting \(y=2(L_0+d)\) and solving
\(y\ee^{-y}=4\Delta\ee^{-2d}/c_0\).

\section{The moving fluctuation and extreme-action scales}
\label{sec:scales}

\begin{lemma}[Existence, size, and exact scale identity]
\label{lem:scales}
For every fixed \(K\), all sufficiently large \(n\), and
\(|\eps|\log n\le K\), the equation
\(\ee^{2H}=nc_{\eps}\FF(H)\) has a unique root in
\([\tfrac25\log n,\tfrac35\log n]\). At this root,
\begin{equation}
 H=\tfrac12\log n+\tfrac12\log\log n+O_K(1),\qquad
 F\asymp_K H,\qquad r\asymp_K H^{-1}.
 \label{eq:H-size}
\end{equation}
With \(\alpha=2-\eps\),
\begin{equation}
 \left(\frac NB\right)^{\alpha}=r,
 \qquad \frac NB=r^{1/\alpha}\asymp_K H^{-1/2}.
 \label{eq:N-B-identity}
\end{equation}
\end{lemma}
\begin{proof}
On the displayed interval, \(\FF(H)\asymp_K\log n\). The function
\(2H-\log(nc_{\eps})-\log\FF(H)\) changes sign at its endpoints.
Its derivative is \(2-\ee^{\eps H}/\FF(H)>1\) for large \(n\).
Taking logarithms of the defining equation gives \eqref{eq:H-size}.
Finally,
\[
 (N/B)^{\alpha}=\frac{nc_{\eps}}{B^{2-\eps}}
 =\frac{B^{\eps}}F=r.
\]
Since \(\eps=O_K(1/H)\),
\(r^{1/\alpha}/\sqrt r=\exp(O_K(\log H/H))\), proving the comparison.
\end{proof}

\begin{corollary}[Explicit crossover function]
\label{cor:explicit-scales}
Write \(\theta_n=\eps\log n\). Uniformly for \(|\theta_n|\le K\),
\begin{equation}
 B^2=\frac{c_{\eps}n\log n}{2}\hc(\theta_n/2)
 \left[1+O_K\left(\frac{\log\log n}{\log n}\right)\right].
 \label{eq:B-explicit}
\end{equation}
If \(\theta_n\to\theta\), then
\begin{equation}
 \frac N{\sqrt n}\longrightarrow \sqrt{c_0}\ee^{\theta/4}.
 \label{eq:N-explicit}
\end{equation}
\end{corollary}
\begin{proof}
Insert \(H=\tfrac12\log n+O_K(\log\log n)\) into
\(F=H\hc(\eps H)\). The function \(\hc\) and its derivatives are bounded
on the relevant compact interval. For the second formula, take logarithms:
\[
 \log(N/\sqrt n)=\frac{\log c_{\eps}}{2-\eps}
       +\frac{\eps\log n}{2(2-\eps)}.
\]
Analyticity implies \(c_{\eps}\to c_0\).
\end{proof}

\section{A uniform local limit theorem with its first central correction}
\label{sec:coefficients}

We now prove the needed triangular-array limit directly. Applying a fixed-law
central or stable limit theorem separately for each \(\eps\) would not give
the uniform result.

\begin{lemma}[Uniform Fourier damping]
\label{lem:damping}
Let \(\chi_n(t)=\E\exp(\ii t(S_n-n)/B)\). For large \(n\) in the window,
there exist positive constants \(b_1,b_2\), independent of \(n,\eps\), such
that
\begin{align}
 |\chi_n(t)|&\le\ee^{-b_1t^2},&& |t|\le B^{1/2},\label{eq:damping-small}\\
 |\chi_n(t)|&\le\ee^{-b_2n/B},&& B^{1/2}\le|t|\le\pi B.
 \label{eq:damping-large}
\end{align}
The same estimates hold for
\(\E\exp(\ii t(S_{n,M}-n)/B)\), uniformly when \(M/N\) is in a fixed
compact subinterval of \((0,\infty)\).
\end{lemma}
\begin{proof}
The real logarithm of the modulus is
\(-n\sum w_j(1-\cos(jt/B))\). On \(|y|\le1\),
\(1-\cos y\ge c y^2\). Retain actions
\(J<j\le R\), where \(R=\lfloor B/\max(1,|t|)\rfloor\), or the
minimum of this number and \(M\) in the cutoff case.
For \(|t|\le B^{1/2}\), this gives \(\log R\ge H/2-O(1)\), because
\(M\asymp N\gg B^{1/2}\). Integral comparison implies
\[
 \sum_{J<j\le R}j^{-1+\eps}\ge c_K F.
\]
The retained exponent is therefore at most
\(-c_Knc_{\eps}Ft^2/B^2=-c_Kt^2\).
For the remaining interval, retain only \(j=1\).
Since \(|t/B|\le\pi\),
\(1-\cos(t/B)\ge c t^2/B^2\), and \(w_1\) has a uniform positive lower
bound. The exponent is at most \(-c nt^2/B^2\le-c n/B\).
\end{proof}

\begin{theorem}[Uniform local limit and central coefficient]
\label{thm:coefficient}
For the admissible family and \(|\eps|\log n\le K\),
\begin{equation}
 \sup_{k\in\mathbb Z}
 \left|B\Pp(S_n=n+k)-\ph(k/B)\right|\longrightarrow0,
 \qquad \ph(x)=\frac{\ee^{-x^2/2}}{\sqrt{2\pi}},
 \label{eq:uniform-llt}
\end{equation}
uniformly in the parameter window. At the center the sharper expansion is
\begin{equation}
 u_n=\frac{\rho_{\eps}^{-n}}{\sqrt{2\pi}\,nB}
 \left[1-\frac1F\left\{
 \frac{E(1-\gamma+\log2)}4+\frac{\gamma+v_{\eps}}2
 \right\}+O_K(H^{-2})\right].
 \label{eq:coefficient-main}
\end{equation}
\end{theorem}
\begin{proof}
For \(t\ne0\), use the principal logarithm of \(-\ii t\), and put
\(L_t=\log(-\ii t)\). The exact cumulant is
\[
 \Psi_n(t)=n\sum_{j\ge1}w_j
 (\ee^{\ii tj/B}-1-\ii tj/B).
\]
Lemma~\ref{lem:kernel}, with \(s=-\ii t/B\), gives
\begin{equation}
 \Psi_n(t)=-\frac{t^2}{2}
 +\frac{t^2}{2F}
 \{EL_t-(\tfrac32-\gamma)E-\gamma-v_{\eps}\}
 +\mathcal E_n(t).
 \label{eq:cf-expansion}
\end{equation}
Here we used
\[
 \FF(H-L_t)=F-EL_t+
 O\bigl(|\eps|E|L_t|^2\ee^{|\eps L_t|}\bigr).
\]
On \(0<|t|\le T=H^{1/12}\), the error is bounded by
\begin{equation}
 |\mathcal E_n(t)|\le
 \frac{C_K}{H^2}|t|^2(1+|\log|t||^2)
 \max(|t|^{1/4},|t|^{-1/4})
 +\frac{C_K|t|^3}{BF}.
 \label{eq:cf-error}
\end{equation}
The last term is the analytic remainder. The replacement of
\(g_{\eps},z_{\eps}\) by their endpoint values contributes the first bound.
We may shrink the parameter interval so that \(|\eps|\le1/4\).
The singularity at \(t=0\) in this bound is integrable.

The displayed correction in \eqref{eq:cf-expansion} is uniformly small on
this central interval. Expanding its exponential, and using
Lemma~\ref{lem:damping}, gives an integrated error \(O_K(H^{-2})\).
Indeed, its square is bounded by \(H^{-2}\) times an integrable
Gaussian--power--logarithmic function. The integral of
\eqref{eq:cf-error} has the same bound; \(1/(BF)\) is smaller than every
fixed negative power of \(H\). The part \(|t|>T\) contributes
\(O_K(\ee^{-bT^2})+O(B\ee^{-bn/B})\), again smaller than \(H^{-2}\).
These estimates justify termwise integration, rather than merely a
pointwise characteristic-function limit.

At the center, Fourier inversion is
\[
 B\Pp(S_n=n)=\frac1{2\pi}\int_{-\pi B}^{\pi B}\ee^{\Psi_n(t)}\,dt.
\]
The imaginary part of \(L_t\) is odd and integrates to zero. For a standard
normal random variable \(Z\),
\begin{equation}
 \E(Z^2\log|Z|)=1-\frac\gamma2-\frac{\log2}{2}.
 \label{eq:Gaussian-log-moment}
\end{equation}
This follows by differentiating
\(\E|Z|^p=2^{p/2}\Gamma((p+1)/2)/\sqrt\pi\) at \(p=2\).
Consequently the normalized first correction is
\[
 \frac1{2F}
 \left[E\left(1-\frac\gamma2-\frac{\log2}2\right)
      -(\tfrac32-\gamma)E-\gamma-v_{\eps}\right]
 =-\frac1F\left[\frac{E(1-\gamma+\log2)}4
                 +\frac{\gamma+v_{\eps}}2\right].
\]
Combine this with \eqref{eq:coeff-poisson}.

Finally, the same bounds prove that the characteristic function, extended
by zero outside \([-\pi B,\pi B]\), converges to \(\ee^{-t^2/2}\) in
\(L^1(\R)\), uniformly in the window. Fourier inversion at arbitrary
\(k/B\) then gives \eqref{eq:uniform-llt}, since the exponential Fourier
factor always has modulus one.
\end{proof}

\begin{corollary}[Explicit leading coefficient]
\label{cor:leading-coefficient}
With \(\theta_n=\eps\log n\),
\begin{equation}
 u_n=\frac{\rho_{\eps}^{-n}}
 {n^{3/2}\sqrt{\pi c_{\eps}\log n\,\hc(\theta_n/2)}}
 \left[1+O_K\left(\frac{\log\log n}{\log n}\right)\right].
 \label{eq:leading-coefficient}
\end{equation}
\end{corollary}
The exponential factor \(\rho_{\eps}^{-n}\) must be retained at the actual
parameter. Replacing it by \(\rho_0^{-n}\) is generally invalid even when
\(\eps=O(1/\log n)\), because an \(O(\eps)\) change in \(\log\rho\) is
multiplied by \(n\).

\section{The uniform cutoff correction and cancellation of prefix data}
\label{sec:cutoff}

\begin{lemma}[Deleted mean and retained moments]
\label{lem:cutoff-moments}
Assume \(\ell=M/N\in[\ell_-,\ell_+]\subset(0,\infty)\), and set
\begin{equation}
 b=\log(M/B)=\log\ell+\alpha^{-1}\log r,
 \quad d_M=\frac{n-\E S_{n,M}}B,
 \quad V_M=\frac{\Var S_{n,M}}{B^2}.
 \label{eq:cutoff-vars}
\end{equation}
Uniformly in this range,
\begin{align}
 n\Lambda_M&=\ell^{-\alpha}/\alpha+O_K(M^{-1}),\label{eq:tail-rate}\\
 d_M&=\frac{r^{1/\alpha}\ell^{1-\alpha}}{\alpha-1}+O_K(B^{-1}),
 &d_M^2&=\frac{r}{\ell^2}+O_K(\log H/H^2),\label{eq:deleted-mean}\\
 V_M&=1+\frac{Eb+\gamma+v_{\eps}}F
      +O_K((\log H)^2/H^2).\label{eq:retained-variance}
\end{align}
For \(k=3,4\), the standardized retained cumulants satisfy
\begin{equation}
 \kappa_{k,M}:=\frac n{B^k}\sum_{j\le M}j^kw_j
 =O_K(H^{-k/2}).
 \label{eq:retained-cumulants}
\end{equation}
\end{lemma}
\begin{proof}
For large \(n\), \(M>J\), so the deleted tail is exactly the power tail.
Integral comparison gives, uniformly near \(\alpha=2\),
\[
 \sum_{j>M}j^{-1-\alpha}=M^{-\alpha}/\alpha+O(M^{-1-\alpha}),
 \qquad
 \sum_{j>M}j^{-\alpha}=M^{1-\alpha}/(\alpha-1)+O(M^{-\alpha}).
\]
Use \(nc_{\eps}=N^\alpha\) and \eqref{eq:N-B-identity}.
For the squared deleted mean, the replacements
\[
 r^{2/\alpha}=r\{1+O_K(\log H/H)\},\quad
 \ell^{2(1-\alpha)}=\ell^{-2}\{1+O_K(H^{-1})\},\quad
 (\alpha-1)^{-2}=1+O_K(H^{-1})
\]
give the stated error.

Uniform integral comparison with one endpoint correction gives
\begin{equation}
 \sum_{j\le M}j^{-1+\eps}
 =\FF(\log M)+\gamma+O(|\eps|)+O(M^{-1+\eps}).
 \label{eq:harmonic-uniform}
\end{equation}
For completeness, subtract the integral of \(x^{-1+\eps}\) on \([1,M]\).
The resulting sum of successive sum--integral errors has terms
\(O(j^{-2+\eps})\), and its parameter derivative has terms
\(O(j^{-2+\eps}\log j)\). Both are uniformly summable near zero.
The limiting constant is therefore \(\gamma+O(|\eps|)\); the remaining
endpoint error is \(O(M^{-1+\eps})\). This proves
\eqref{eq:harmonic-uniform} without any interchange of divergent sums.

Add the finite-prefix second moment \(c_{\eps}v_{\eps}\), divide by
\(c_{\eps}F\), and apply \eqref{eq:F-shift} with \(b=O_K(\log H)\).
This proves \eqref{eq:retained-variance}. Finally,
\[
 \frac{nc_{\eps}M^{k-\alpha}}{B^k}
 =\ell^{k-\alpha}r^{k/\alpha}=O_K(H^{-k/2}),\qquad k=3,4.
\]
The finite-prefix contribution \(O(n/B^k)\) is smaller. Integral comparison
for the retained powers completes the proof.
\end{proof}

\begin{lemma}[Central probability for the cutoff]
\label{lem:cutoff-central}
In the same parameter range,
\begin{equation}
 B\sqrt{2\pi}\,\Pp(S_{n,M}=n)
 =1-\frac{Eb+\gamma+v_{\eps}}{2F}
    -\frac{E}{2F\ell^2}
    +O_K((\log H)^2/H^2).
 \label{eq:cutoff-central}
\end{equation}
\end{lemma}
\begin{proof}
The real and imaginary parts of the centered cutoff cumulant obey
\begin{align*}
 \Rea\Psi_{n,M}(t)&=-V_Mt^2/2+O(\kappa_{4,M}t^4),\\
 \Im\Psi_{n,M}(t)&=-d_Mt+O(\kappa_{3,M}|t|^3).
\end{align*}
These inequalities follow directly from the Taylor remainders for cosine
and sine. Fourier inversion at the center integrates
\(\ee^{\Rea\Psi_{n,M}}\cos(\Im\Psi_{n,M})\).
On \(|t|\le H^{1/12}\), put \(a_M=V_M-1\).
Its expansion through the required order is
\[
 \ee^{-t^2/2}\{1-a_Mt^2/2-d_M^2t^2/2\}.
\]
The integrated error is bounded by a fixed constant times
\[
 a_M^2+\kappa_{4,M}+d_M^4
       +|d_M|\kappa_{3,M}+\kappa_{3,M}^2
 =O_K((\log H)^2/H^2).
\]
For the phase estimate, use
\(|\cos(x+y)-\cos x|\le |xy|+y^2/2\).
This is why an apparent cubic error of size \(H^{-3/2}\) does not survive
at the center: it is multiplied by the deleted-mean shift or enters
quadratically. The tail of the Fourier integral is negligible by
Lemma~\ref{lem:damping}. Integrating \(t^2\ee^{-t^2/2}\), then using
Lemma~\ref{lem:cutoff-moments}, proves the result.
\end{proof}

\begin{theorem}[Uniform prefix-independent cutoff correction]
\label{thm:cutoff}
For every admissible family, uniformly for \(|\eps|\log n\le K\) and
\(\ell=M/N\in[\ell_-,\ell_+]\subset(0,\infty)\),
\begin{equation}
 \frac{u_{n,M}}{u_n}
 =\ee^{-\ell^{-\alpha}/\alpha}
 \left[
 1+\frac r4\left\{\log\frac{2}{r\ell^2}+1-\gamma-\frac2{\ell^2}\right\}
 +O_K\left(\frac{(\log H)^2}{H^2}\right)
 \right].
 \label{eq:cutoff-theorem}
\end{equation}
In particular, the displayed normalized correction is the same for all
finite-prefix changes, after using that family's actual \(c_{\eps}\),
\(N\), and \(B\).
\end{theorem}
\begin{proof}
Write the full central probability as
\((B\sqrt{2\pi})^{-1}(1+C/F+O(H^{-2}))\), where
\[
 C=-\frac{E(1-\gamma+\log2)}4-\frac{\gamma+v_{\eps}}2.
\]
Divide \eqref{eq:cutoff-central} by this expression. The numerator's
\(-(\gamma+v_{\eps})/(2F)\) cancels the identical denominator term.
What remains is
\begin{equation}
 \frac{\Pp(S_{n,M}=n)}{\Pp(S_n=n)}
 =1+\frac EF\left\{-\frac b2
        +\frac{1-\gamma+\log2}{4}-\frac1{2\ell^2}\right\}
   +O_K((\log H)^2/H^2).
 \label{eq:probability-ratio}
\end{equation}
Now \(b=\log\ell+\alpha^{-1}\log r\), and
\(\alpha^{-1}=1/2+O_K(H^{-1})\). Multiplication by \(r=O_K(H^{-1})\)
shows that replacing \(\alpha^{-1}\log r\) by \(\tfrac12\log r\)
costs only \(O_K(\log H/H^2)\). Therefore the displayed term in
\eqref{eq:probability-ratio} becomes
\[
 \frac r4\left\{\log\frac{2}{r\ell^2}+1-\gamma-\frac2{\ell^2}\right\}.
\]
Lastly, \eqref{eq:tail-rate} replaces the survival factor in
\eqref{eq:cutoff-poisson} by \(\ee^{-\ell^{-\alpha}/\alpha}\) with relative
error \(O_K(M^{-1})\), smaller than the stated remainder.
\end{proof}

\begin{remark}[The exponent in the survival factor must not be frozen]
Although \(\alpha\to2\), replacing
\(\ee^{-\ell^{-\alpha}/\alpha}\) by \(\ee^{-1/(2\ell^2)}\) can change a
term of order \(1/H\). Such a replacement is allowed for the limiting
law, but not generally for the first correction. Formula
\eqref{eq:cutoff-theorem} keeps this contribution exactly.
\end{remark}

\begin{corollary}[Lossless truncation and the raw square-root scale]
\label{cor:lossless}
In the bounded window,
\begin{equation}
 \frac{u_{n,M_n}}{u_n}\longrightarrow1
 \quad\Longleftrightarrow\quad \frac{M_n}{N_n}\longrightarrow\infty.
 \label{eq:iff-lossless}
\end{equation}
If \(M_n/N_n\to0\), the ratio tends to zero. If \(M_n/N_n\to\ell>0\),
the limit is \(\ee^{-1/(2\ell^2)}\). Equivalently, when
\(\eps\log n\to\theta\) and \(M_n/\sqrt n\to m>0\),
\begin{equation}
 \frac{u_{n,M_n}}{u_n}\longrightarrow
 \exp\left(-\frac{c_0\ee^{\theta/2}}{2m^2}\right).
 \label{eq:raw-cutoff}
\end{equation}
\end{corollary}
\begin{proof}
The coefficient ratio is increasing in \(M\) and lies in \([0,1]\).
For each fixed \(L>0\), Theorem~\ref{thm:cutoff} applies at
\(M=\lfloor LN\rfloor\). If \(M_n/N_n\to\infty\), use this cutoff as
a lower bound and then let \(L\to\infty\). Conversely, a bounded
subsequence of \(M_n/N_n\) is bounded above by some finite \(L\), giving
an upper limit at most \(\ee^{-1/(2L^2)}<1\).
The zero-limit statement follows by an upper bound at each fixed small
\(L\) and then \(L\downarrow0\). The remaining assertions follow directly
from the theorem and \eqref{eq:N-explicit}.
\end{proof}

\section{A two-term minimal action budget}
\label{sec:budget}

Fix a target retained fraction \(\tau\in(0,1)\), independent of \(n\).
Since actions above \(n\) cannot contribute to \([q^n]\), there is always
an integer \(M\le n\) retaining the whole coefficient. Thus define
\begin{equation}
 M_{\min}(n,\eps;\tau)=\min\{M\ge1:u_{n,M}/u_n\ge\tau\}.
 \label{eq:min-budget}
\end{equation}

\begin{corollary}[Uniform two-term budget]
\label{cor:budget}
Let \(L_\tau=-\log\tau\) and
\(\ell_0=(\alpha L_\tau)^{-1/\alpha}\). In the bounded parameter window,
\begin{align}
 \frac{M_{\min}}N
 =\ell_0\biggl[1-&\frac{r}{4\alpha L_\tau}
 \left\{\log\frac{2}{r\ell_0^2}+1-\gamma-\frac2{\ell_0^2}\right\}
 \nonumber\\
 &+O_{K,\tau}\left(\frac{(\log H)^2}{H^2}\right)\biggr].
 \label{eq:budget-formula}
\end{align}
\end{corollary}
\begin{proof}
The leading cutoff law brackets \(M_{\min}/N\) inside a fixed compact
interval around \(\ell_0\). The logarithm of
\eqref{eq:cutoff-theorem} is
\[
 -\ell^{-\alpha}/\alpha
 +\frac r4 h(\ell,r)+O_K((\log H)^2/H^2),\qquad
 h(\ell,r)=\log\frac{2}{r\ell^2}+1-\gamma-\frac2{\ell^2}.
\]
The derivative of the leading term at \(\ell_0\) is
\(\ell_0^{-\alpha-1}>0\). The first displacement of its crossing of
\(\log\tau\) is therefore
\(-r\ell_0^{\alpha+1}h(\ell_0,r)/4\).
Taylor's formula on that compact interval bounds the remaining displacement
by \(O_{K,\tau}((\log H)^2/H^2)\). Monotonicity brackets the integer
crossing between the two continuous error bounds; the rounding error is
at most \(1/N\), which is smaller. Since
\(\ell_0^\alpha=(\alpha L_\tau)^{-1}\), this is
\eqref{eq:budget-formula}.
\end{proof}

This is an asymptotic budget, not a finite-\(n\) interval certificate.
For an actual finite computation, a positive coefficient recurrence gives a
separate way of testing the requested retained fraction. Section~\ref{sec:checks}
describes such a recurrence. A uniform numerical value for the constants in
the asymptotic error has not been certified here.

\section{Conditional large actions converge to a Poisson process}
\label{sec:extremes}

The cutoff theorem is a largest-action statement. The local limit theorem
allows an extension to all macroscopic extreme actions at once.

\begin{theorem}[Conditional Poisson limit]
\label{thm:PPP}
Let \(\eps=\eps_n\) with \(|\eps_n|\log n\le K\). Conditional on
\(S_n=n\), define the random counting measure
\begin{equation}
 \Xi_n=\sum_{j\ge1}Z_j\,\delta_{j/N}
 \quad\text{on }(0,\infty).
 \label{eq:point-process}
\end{equation}
Then \(\Xi_n\) converges in distribution, for the vague topology on locally
finite counting measures, to a Poisson point process with intensity
\begin{equation}
 \nu(dx)=x^{-3}\,dx.
 \label{eq:PPP-intensity}
\end{equation}
In particular, if \(J_{(k)}\) is the \(k\)-th largest action, counting
multiplicities and padded with zeros, then for each \(\ell>0\),
\begin{equation}
 \Pp(J_{(k)}/N\le\ell\mid S_n=n)
 \longrightarrow
 \ee^{-1/(2\ell^2)}
 \sum_{m=0}^{k-1}\frac{(1/(2\ell^2))^m}{m!}.
 \label{eq:kth-largest}
\end{equation}
\end{theorem}
\begin{proof}
For finitely many distinct indices and integers \(m_j\ge0\), the Poisson
factorial-moment identity is exact:
\begin{equation}
 \E\left[\prod_j(Z_j)_{m_j}\,\middle|\,S_n=n\right]
 =\prod_j(nw_j)^{m_j}
 \frac{\Pp(S_n=n-\sum_jjm_j)}{\Pp(S_n=n)}.
 \label{eq:conditional-factorial}
\end{equation}
It follows by shifting each count in its Poisson probability mass function;
\((z)_m=z(z-1)\cdots(z-m+1)\).

Consider disjoint intervals \(I_1,\dots,I_d\) whose closures lie in
\((0,\infty)\). A fixed mixed factorial moment of their counts has shifts
of size \(O(N)\). Since \(N/B\to0\),
Theorem~\ref{thm:coefficient} gives
\[
 \frac{\Pp(S_n=n-k)}{\Pp(S_n=n)}\longrightarrow1
 \quad\text{uniformly for } |k|\le C N
\]
for every fixed \(C\). Summing \eqref{eq:conditional-factorial} over the
indices in the intervals yields the mixed factorial moments of independent
Poisson variables, because
\begin{equation}
 n\sum_{j/N\in I_a}w_j
 =\frac1N\sum_{j/N\in I_a}(j/N)^{-1-\alpha}
 \longrightarrow\int_{I_a}x^{-3}\,dx.
 \label{eq:intensity-riemann}
\end{equation}
The finite prefix is eventually outside every such interval.

There is no moment-determinacy gap: the uniform local limit also bounds
\(\sup_k\Pp(S_n=k)/\Pp(S_n=n)\) by a fixed constant. Thus every mixed
factorial moment is bounded by that constant times the corresponding
product of the unconditioned means to the required powers. The exponential
series for the joint count generating function is consequently dominated
on a neighborhood of one. Its limit is the independent-Poisson generating
function. The same first-moment bound gives tightness of the number of
points in each compact set. Approximation of a nonnegative continuous
compactly supported test function by interval step functions proves
convergence of the Laplace functionals, and hence the asserted vague limit.

To pass from compact intervals to the entire tail \((\ell,\infty)\), note
from the same identity and probability-ratio bound that
\[
 \E[\Xi_n((L,\infty))\mid S_n=n]
 \le Cn\Lambda_{\lfloor LN\rfloor}\le C_K L^{-7/4}
\]
for \(L\ge1\) and all sufficiently large \(n\).
The omitted upper tail is therefore negligible as \(L\to\infty\).
The limiting number of actions above \(\ell\) is Poisson with mean
\(\int_\ell^\infty x^{-3}dx=1/(2\ell^2)\). At most \(k-1\) such actions
is exactly the event in \eqref{eq:kth-largest}.
\end{proof}

The conditioning is essential in the formulation: coefficients are not
unconditioned probabilities. The proof shows why it becomes invisible at
the extreme-action scale in this window. Any fixed number of extreme
actions shifts the total by \(O(N)=o(B)\), too little to alter the leading
local Gaussian density. The first cutoff correction measures part of the
remaining, smaller interaction.

\section{Interpretation and the finite-variance-side surprise}
\label{sec:interpretation}

\subsection{Three different limits}
For fixed \(1<\alpha<2\), the classical coefficient mechanism is stable,
with a power scale \(n^{1/\alpha}\). At exactly \(\alpha=2\), the central
scale is \(\sqrt{n\log n}\). In the window of this article,
\(\alpha=2-\theta_n/\log n\), the central limit is Gaussian but its scale
retains the nontrivial factor \(\hc(\theta_n/2)\). These are three different
asymptotic assertions. A bounded-window theorem does not by itself justify
letting \(\theta_n\to+\infty\) and claiming recovery of a fixed-index
stable law.

The action scale is different again:
\(N\sim\sqrt{c_0n}\ee^{\theta/4}\), whereas
\(B^2\sim(c_0n\log n/2)\hc(\theta/2)\).
Thus a cutoff much larger than \(\sqrt n\), in the normalized sense of
Corollary~\ref{cor:lossless}, preserves the coefficient even when it is much
smaller than the central fluctuation scale. Confusing a largest-jump scale
with a sum-fluctuation scale overestimates the required action budget.

\subsection{Finite variance need not give the correct central scale}
Take \(\eps_n=-\kappa/\log n\), with fixed \(\kappa>0\). Every member of
this triangular family has \(\alpha_n>2\) and finite variance. Nevertheless,
its total variance is not asymptotic to \(B^2\).
Indeed,
\begin{equation}
 \Var S_n=nc_{\eps}\{\zeta(1-\eps)+v_{\eps}\},\qquad
 \frac{\Var S_n}{B^2}\longrightarrow
 \frac1{1-\ee^{-\kappa/2}}.
 \label{eq:variance-mismatch}
\end{equation}
This follows from the zeta pole and \eqref{eq:B-explicit}.
Theorem~\ref{thm:coefficient}, together with its characteristic-function
proof, implies \((S_n-n)/B\Rightarrow N(0,1)\). Normalizing instead by the
actual standard deviation therefore gives a Gaussian with limiting variance
\(1-\ee^{-\kappa/2}\), not one.

There is no contradiction. Rare actions far beyond the central scale retain
a nonvanishing share of the second moment, although their probabilities
vanish. For example, for fixed \(a>0\), integral comparison gives
\[
 \frac{nc_{\eps}\sum_{j>aB}j^{-1+\eps}}{\Var S_n}
 \longrightarrow\ee^{-\kappa/2}.
\]
The numerator is the contribution of those jumps to the Poisson variance.
This is an explicit failure of uniform integrability of the normalized
second moments. A fixed-distribution finite-variance intuition is therefore
not an adequate substitute for the triangular Fourier proof.

\subsection{Why finite-prefix data cancel only in the displayed ratio}
The inverse expansion contains \(v_{\eps}\), and the central coefficient
contains \(v_{\eps}\). Neither is prefix-independent. The ratio is special:
the same retained finite-prefix variance enters its numerator and denominator
with the same coefficient, and disappears upon division. Prefix changes also
alter \(c_{\eps}\), \(\rho_{\eps}\), \(B\), and \(N\). The theorem never
claims invariance when these normalizations are held artificially fixed.
Nor does it assert that prefix data cancel to every higher order.

\section{Reproducible checks and numerical diagnostics}
\label{sec:checks}

\subsection{What was actually checked}
The supplied program \texttt{code/verify.py} executed 17 exact symbolic checks:
14 comparisons between direct formal iteration of the feedback equation and
the Lagrange coefficient identity, plus three checks of the Gaussian constant
and the prefix-cancellation algebra. The formal comparisons use two rational
finite-action families and degrees one through seven. They test finite
identities; they do not certify the analytic theorems.

For an independent finite-\(n\) diagnostic, if
\(p_k=\Pp(S_n=k)\), differentiation of the probability generating function
gives the positive recurrence
\begin{equation}
 p_0=\ee^{-nG(1)},\qquad
 p_k=\frac nk\sum_{j=1}^k jw_jp_{k-j}.
 \label{eq:prob-recurrence}
\end{equation}
The cutoff version restricts the sum to \(j\le M\) and changes \(p_0\)
accordingly. The program ran 18 such diagnostics, for
\(n=512,2048,4096\), \(\theta=\eps\log n\in\{-1,0,1\}\), and
\(\eta\in\{0,2\}\) in \eqref{eq:test-family}. In six cross-check cases at
\(n=512\), recurrence and independently evaluated central Fourier integrals
agreed with maximum absolute discrepancy
\(3.3306690738754696\times10^{-16}\), across the tested ratios and normalized
central probabilities.

\subsection{Large-order cutoff checks}
The larger orders in Table~\ref{tab:large} are evaluated by a scaled Fourier
integral, not by constructing \(10^{32}\) coefficients. The full cumulant uses
the desingularized polylogarithm formula and 16 analytic terms. The cutoff
cumulant uses 60 terms, with finite moments evaluated by direct sums or
14-term Euler--Maclaurin evaluation. Numerical integration is over
\(0\le t\le14\). These implementation truncations and floating-point
quadrature estimates are not rigorous enclosures.

All displayed cutoff runs take \(M\) to be the nearest integer to \(N\),
but evaluate both approximations at the actual \(\ell=M/N\). The column
``leading'' keeps \(\alpha\) in the survival exponential. ``Corrected''
uses the full displayed expression in Theorem~\ref{thm:cutoff}.

\begin{table}[htbp]
\centering\small
\begin{tabular}{rrrrr}
\toprule
$n$ & $\theta$ & Computed ratio & Leading & Corrected\\
\midrule
$10^{8}$ & -2 & 0.639884263 & 0.622371360 & 0.639837832\\
$10^{8}$ & 0 & 0.628290097 & 0.606533148 & 0.627967224\\
$10^{8}$ & 2 & 0.610090576 & 0.589362780 & 0.611719991\\
$10^{16}$ & -2 & 0.626485822 & 0.614597962 & 0.626437090\\
$10^{16}$ & 0 & 0.622905481 & 0.606530659 & 0.622684654\\
$10^{16}$ & 2 & 0.618154572 & 0.598128083 & 0.617962252\\
$10^{32}$ & -2 & 0.618213163 & 0.610604999 & 0.618168907\\
$10^{32}$ & 0 & 0.617565238 & 0.606530660 & 0.617441970\\
$10^{32}$ & 2 & 0.617011900 & 0.602372524 & 0.616782127\\
\bottomrule
\end{tabular}
\caption{Uniform-window cutoff diagnostics for the pure family. These are
floating-point evaluations, not certified exact values.}
\label{tab:large}
\end{table}

The convergence is logarithmically slow. A very large value of \(n\) need not
make the uncorrected cutoff law highly accurate. For example, at
\(n=10^{16}\), \(\theta=2\), the computed ratio is approximately
\(0.618154572\), compared with \(0.598128083\) from the leading term and
\(0.617962252\) from the correction. This is evidence consistent with the
formula, not a substitute for its uniform error proof.

\subsection{Inverse checks}
For the inverse diagnostics, \(D=\log(1/\Delta)\) and
\(\theta=\eps D\). A high-precision root solve in the variable
\(L=\log(1/s)\), using the desingularized expression, was compared with
both inverse approximations. The analytic remainder was evaluated through
order 16. Table~\ref{tab:inverse} gives signed relative errors, approximation
divided by the computed reference value minus one.

\begin{table}[htbp]
\centering\small
\begin{tabular}{rrrr}
\toprule
$D$ & $\theta$ & Leading relative error & Corrected relative error\\
\midrule
20 & -2 & 0.07631914 & -0.01277858\\
20 & 0 & 0.10242388 & -0.01089807\\
20 & 2 & 0.17846262 & -0.00151471\\
40 & -2 & 0.04410638 & -0.00418078\\
40 & 0 & 0.05972766 & -0.00392070\\
40 & 2 & 0.09761422 & -0.00194281\\
80 & -2 & 0.02490134 & -0.00133200\\
80 & 0 & 0.03419827 & -0.00133957\\
80 & 2 & 0.05423873 & -0.00095531\\
\bottomrule
\end{tabular}
\caption{Two-term inverse diagnostics for $\eta=0$. The joint parameter here
is $\theta=\varepsilon D$, not $\varepsilon\log n$.}
\label{tab:inverse}
\end{table}

\subsection{Reproduction and limitations}
The recorded environment was Python 3.13.5, NumPy 2.3.5, mpmath 1.3.0,
SciPy 1.17.0, and SymPy 1.14.0. Constants and finite moments were evaluated
with 75 mpmath decimal digits; the recurrence used NumPy long double;
quadrature used ordinary double precision. The complete CSV data, exact-check
list, environment record, and run summary are included. The default
recurrence order is deliberately limited to avoid underflow of its initial
Poisson mass on the recorded platform. The script reports underflow rather
than silently returning a spurious answer.

The numerical program is not a Lean proof, a certified arbitrary-precision
interval implementation, or a proof of the asymptotic error constants.
Theorems~\ref{thm:inverse}, \ref{thm:chart}, \ref{thm:coefficient},
\ref{thm:cutoff}, and \ref{thm:PPP} stand on their written proofs.

\section{Proof dependencies, repository scope, and formalization}
\label{sec:audit}

\subsection{Dependency structure}
The inverse theory depends on the exact polylogarithm expansion, critical
normalization, elementary uniform estimates, and the analytic
implicit-function theorem. The coefficient theory depends on the Lagrange
identity, the same desingularized expansion, Fourier inversion, and the
explicit damping lemma. The cutoff correction then uses only retained
moments and division of the two local expansions. The minimal budget uses
monotonicity. The point-process result uses the uniform local limit and
an exact Poisson factorial-moment identity.

No result from the companion marginal manuscript is used as a proof premise.
Its open question is the research target, and its displayed endpoint formula
is an overlap check. The argument here rederives that endpoint while proving
uniformity in the moving parameter. No unverified theorem from the full
canonical \repo\ volume is needed.

\subsection{Scope of the repository inspection}
The inspected public snapshot is commit
\texttt{0c973d8f5e7ef4550f4df1bf486d890037fd9f14}.
The inspection covered the transseries project README, the opening portion
of the group index, the relevant package inventory, and the complete README
for the fixed-stable critical-Hahn package. The latter identifies its theorem
scope and proof boundaries. The private companion
\texttt{Marginal\_Critical\_Transseries.tex} was inspected through its abstract,
main-result summary, and explicit further-research question. We did not audit
the entire canonical volume or all source files in the repository. No
repository files were modified.

This separation matters. An entry in the Lean crosswalk is evidence about
that entry, not automatic formal verification of a neighboring analytic
research claim. The present results are not represented as already available
in the repository's Lean development.

\subsection{A realistic formalization route}
A useful Lean development can be split into four interfaces.
First, formalize the finite coefficient identity and positivity using formal
power series and a truncated exponential. Second, formalize the elementary
scale algebra, including \((N/B)^\alpha=r\), and the cancellation of the
finite-prefix moment. Third, supply analytic interfaces for the removable
polylogarithm kernel and the parameter-uniform implicit theorem. Fourth,
formalize a reusable lattice Fourier lemma with integrable remainder and
damping hypotheses, and instantiate it for the full and cutoff sums.

The hard analytic obligations should remain explicit: normal convergence
in the parameter, a uniform Fourier tail, the integrated central remainder,
and the probability-ratio bound used in conditional factorial moments.
Verifying a symbolic coefficient identity does not discharge any of these.
The supplied exact checks suggest small initial formalization targets, not
a claim of an existing machine-checked asymptotic proof.

\section{Further research questions}
\label{sec:further}

\paragraph{1. Leave the bounded window and match the stable law.}
Extend the argument to \(\eps\log n\to+\infty\) while \(\eps\to0\),
with explicit uniform errors. Determine the maximal range in which a
Gaussian approximation remains accurate after the deformed normalization,
and construct an overlap expansion with the fixed-index stable density.
The constants hidden in the bounded-window estimates cannot simply be
assumed uniform as the window grows.

\paragraph{2. Compute the next correction and locate prefix reentry.}
Expand the full Fourier cumulant to order \(H^{-2}\), retaining the parameter
variation of the regularized gamma and zeta constants. Combine it with the
fourth retained cumulant and the squared deleted-mean shift. Determine exactly
which finite-prefix moments survive at that order. The cancellation proved
here gives a nontrivial check on any such calculation.

\paragraph{3. Add a slowly varying tail.}
Replace the eventually exact tail by
\(c j^{-3+\eps}L(j)\), with quantified uniform regularity of \(L\).
Identify necessary and sufficient conditions for the ratio between the
extreme-action and central scales to vanish. Establish which slowly varying
perturbations preserve a finite-generator inverse chart and which require
new core inverses. A pointwise tail asymptotic alone is unlikely to control
the first correction.

\paragraph{4. Allow a drifting distance from criticality.}
Replace \(\sum jw_j=1\) by a mean differing from one on the \(B/n\) scale.
The leading local Gaussian is then sampled away from zero. Determine a joint
coupling--index--cutoff expansion, including whether the odd Fourier term
creates a prefix-dependent correction before order \(H^{-2}\).
The cancellation of the cubic contribution at the center must be revisited.

\paragraph{5. Quantify conditional Poisson convergence.}
The point-process theorem is qualitative. Derive explicit total-variation
or bounded-test-function rates for counts on intervals separated from zero,
and compute the first conditioning-induced covariance between two disjoint
large-action windows. The exact factorial-moment formula and the shifted
local Fourier expansion provide a direct route.

\paragraph{6. Develop complex continuation of the deformed core.}
The correction germ is holomorphic, but the core is selected here on its
large real branch. Classify its nearby complex branches on a logarithmic
cover, locate parameter-dependent turning points, and compare the monodromy
with the Lambert endpoint and nonintegral Hahn regimes. A global two-sheet
or Stokes statement would require additional work; it is not implicit in
Theorem~\ref{thm:chart}.

\paragraph{7. Replace asymptotic budgets by certified finite budgets.}
Make every constant in the damping and Taylor estimates explicit, and combine
them with interval arithmetic for the positive recurrence or Fourier
integrals. The goal is a proof-producing routine returning a guaranteed
smallest or near-smallest \(M\) for a prescribed coefficient fraction. The
current two-term formula is a principled initial estimate, not that certificate.

\paragraph{8. Signed and multi-type feedback.}
For signed weights the Poisson model and coefficient monotonicity disappear;
for multiple feedback variables the criticality condition becomes a spectral
condition. Determine which pole-cancellation and inverse-chart arguments
survive, and identify replacement hypotheses for the scalar positive Fourier
proof. These are distinct extensions, not consequences of analytic
continuation of the positive-weight formulas.

\section{Conclusion}
The joint stable--Gaussian endpoint question has an explicit answer for
positive finite-prefix power-tail feedback. The correct shared kernel is
\((\ee^{\eps L}-1)/\eps\). It simultaneously controls critical inversion and
central coefficient growth. The associated extreme-action scale is smaller
than the fluctuation scale by a logarithmic factor. At that smaller scale,
the first cutoff correction is universal after normalization: finite-prefix
second moments cancel, even though they remain visible in the two separate
local expansions. The same separation of scales makes the conditioned large
actions asymptotically Poisson.

These conclusions resolve the specified bounded-window follow-up, including
its first interpolating cutoff correction, within the stated family and with
conventional proofs. They do not assert an unrestricted transseries theorem,
a global priority result, or a formal verification that has not been carried
out. The exact formulas and explicit proof obligations make the next
extensions concrete.

\appendix
\section{A geometric truncation bound for the convergent chart}
\label{app:chart}
The analyticity statement in Theorem~\ref{thm:chart} implies an ordinary
geometric error bound, distinct from the large-order asymptotic errors.
After shrinking its domain, there exist constants \(R_0,C>0\) such that,
uniformly for \(|\eps|\le\delta\),
\[
 V(\eps,X,Y,Z)=\sum_{a,b,c\ge0}v_{abc}(\eps)X^aY^bZ^c,
 \qquad |v_{abc}(\eps)|\le C R_0^{-a-b-c}.
\]
Let \(V_{\le m}\) be its truncation to total degree at most \(m\), and put
\(t=\max(|X|,|Y|,|Z|)/R_0<1\). Cauchy's estimate gives
\begin{equation}
 |V-V_{\le m}|\le
 C\sum_{k>m}\binom{k+2}{2}t^k
 \le C\frac{(m+3)^2t^{m+1}}{(1-t)^3}.
 \label{eq:geometric-tail}
\end{equation}
The second inequality follows by writing \(k=m+1+l\), bounding
\(\binom{m+l+3}{2}\) by \((m+3)^2(l+1)^2/2\), and summing
\(\sum_{l\ge0}(l+1)^2t^l=(1+t)/(1-t)^3\).
A bound on the relative error of \(s_0\ee^{-V_{\le m}}\) follows from
\(|\ee^z-1|\le |z|\ee^{|z|}\). This is a convergent analytic truncation
statement. It does not turn the \(H^{-1}\) coefficient expansion into a
convergent series, nor supply numerical values of \(C,R_0\) without a
quantitative implicit-function analysis.

\phantomsection
\addcontentsline{toc}{section}{References}
\begin{thebibliography}{9}
\bibitem{repo-readme}
Vladimir Reshetnikov, \emph{ProveIt: Transseries project README}, inspected
29 September 2026, snapshot
\texttt{0c973d8f5e7ef4550f4df1bf486d890037fd9f14}.
\url{https://github.com/VladimirReshetnikov/ProveIt/blob/0c973d8f5e7ef4550f4df1bf486d890037fd9f14/Analysis/Transseries/README.md}.

\bibitem{repo-index}
Vladimir Reshetnikov, \emph{ProveIt: Series and transseries group index},
same snapshot; opening 140 source lines and relevant package inventory
inspected.
\url{https://github.com/VladimirReshetnikov/ProveIt/blob/0c973d8f5e7ef4550f4df1bf486d890037fd9f14/Analysis/Transseries/docs/series-and-transseries/README.md}.

\bibitem{repo-stable}
\emph{Beyond Finite-Action Folds: Critical Hahn Transseries, Stable Sector
Asymptotics, and Sharp Action Budgets}, research package prepared for Vladimir
Reshetnikov, 29 September 2026. Package README inspected at the snapshot above;
used for scope comparison, not as a proof premise.
\url{https://github.com/VladimirReshetnikov/ProveIt/tree/0c973d8f5e7ef4550f4df1bf486d890037fd9f14/Analysis/Transseries/docs/series-and-transseries/Critical_Hahn_Transseries_Beyond_Finite_Action_Folds}.

\bibitem{marginal}
\emph{Marginal Critical Transseries: Convergent Lambert--$W$ Charts, Universal
Cutoff Corrections, and Certified Action Budgets}, unpublished companion draft
prepared for Vladimir Reshetnikov, 29 September 2026. Inspected in the user's
file library as \texttt{Marginal\_Critical\_Transseries.tex}, especially the
main-result summary and the first further-research question. No public
publication or independent review is asserted.
% ed. (2026-09-29): filed location added.
Editorial note (ProveIt, 2026-09-29): now filed (batch~48) in the group of
\cite{repo-index} as the package
\path{Marginal_Critical_Transseries_Action_Budgets/}; the question answered
here is its Question~1.

% ed. (2026-09-29): the two entries below are editorial additions for the editorial note in Section 1.1.
\bibitem{ed:lce}
[Editorial addition, ProveIt, 2026-09-29.]
\emph{The Logarithmic Critical Endpoint: Convergent Lambert Charts,
All-Order Sector Laws, and Sharp Action-Cutoff Corrections}, research
package filed in batch~48 in the group of \cite{repo-index} as
\path{Logarithmic_Critical_Endpoint_Lambert_Charts/}.

\bibitem{ed:cct}
[Editorial addition, ProveIt, 2026-09-29.]
\emph{Confluent Critical Transseries Across the Exponent-Two Boundary},
research package filed in batch~49 in the group of \cite{repo-index} as
\path{Confluent_Critical_Transseries_Exponent_Two_Boundary/}.

\bibitem{dlmf-polylog}
F. W. J. Olver et al. (eds.), \emph{NIST Digital Library of Mathematical
Functions}, Section 25.12, especially Eq. 25.12.12, polylogarithm expansion.
\url{https://dlmf.nist.gov/25.12}. Accessed 29 September 2026.

\bibitem{janson}
Svante Janson, Simply generated trees, conditioned Galton--Watson trees,
random allocations and condensation, \emph{Probability Surveys} \textbf{9}
(2012), 103--252. DOI: \texttt{10.1214/11-PS188}.
\url{https://arxiv.org/abs/1112.0510}. Theorem 17.14, its local-limit proof,
and Example 18.29 provide relevant fixed-exponent and endpoint background.

\bibitem{chakrabarty}
Arijit Chakrabarty and Gennady Samorodnitsky, Understanding heavy tails in a
bounded world or, is a truncated heavy tail heavy or not?,
\emph{Stochastic Models} \textbf{28} (2012), no.~1, 109--143.
DOI: \texttt{10.1080/15326349.2012.646551}.
\url{https://arxiv.org/abs/1001.3218}. Used for historical context on truncation
regimes, not to infer the moving-index corrections proved here.
\end{thebibliography}
\end{document}
